\PassOptionsToPackage{unicode=true}{hyperref} % options for packages loaded elsewhere
\PassOptionsToPackage{hyphens}{url}
%
\documentclass[]{book}
\usepackage{lmodern}
\usepackage{amssymb,amsmath}
\usepackage{ifxetex,ifluatex}
\usepackage{fixltx2e} % provides \textsubscript
\ifnum 0\ifxetex 1\fi\ifluatex 1\fi=0 % if pdftex
  \usepackage[T1]{fontenc}
  \usepackage[utf8]{inputenc}
  \usepackage{textcomp} % provides euro and other symbols
\else % if luatex or xelatex
  \usepackage{unicode-math}
  \defaultfontfeatures{Ligatures=TeX,Scale=MatchLowercase}
\fi
% use upquote if available, for straight quotes in verbatim environments
\IfFileExists{upquote.sty}{\usepackage{upquote}}{}
% use microtype if available
\IfFileExists{microtype.sty}{%
\usepackage[]{microtype}
\UseMicrotypeSet[protrusion]{basicmath} % disable protrusion for tt fonts
}{}
\IfFileExists{parskip.sty}{%
\usepackage{parskip}
}{% else
\setlength{\parindent}{0pt}
\setlength{\parskip}{6pt plus 2pt minus 1pt}
}
\usepackage{hyperref}
\hypersetup{
            pdftitle={A Student's History of Education},
            pdfauthor={Frank Pierrepont Graves},
            pdfborder={0 0 0},
            breaklinks=true}
\urlstyle{same}  % don't use monospace font for urls
\usepackage{longtable,booktabs}
% Fix footnotes in tables (requires footnote package)
\IfFileExists{footnote.sty}{\usepackage{footnote}\makesavenoteenv{longtable}}{}
\usepackage{graphicx,grffile}
\makeatletter
\def\maxwidth{\ifdim\Gin@nat@width>\linewidth\linewidth\else\Gin@nat@width\fi}
\def\maxheight{\ifdim\Gin@nat@height>\textheight\textheight\else\Gin@nat@height\fi}
\makeatother
% Scale images if necessary, so that they will not overflow the page
% margins by default, and it is still possible to overwrite the defaults
% using explicit options in \includegraphics[width, height, ...]{}
\setkeys{Gin}{width=\maxwidth,height=\maxheight,keepaspectratio}
\setlength{\emergencystretch}{3em}  % prevent overfull lines
\providecommand{\tightlist}{%
  \setlength{\itemsep}{0pt}\setlength{\parskip}{0pt}}
\setcounter{secnumdepth}{5}
% Redefines (sub)paragraphs to behave more like sections
\ifx\paragraph\undefined\else
\let\oldparagraph\paragraph
\renewcommand{\paragraph}[1]{\oldparagraph{#1}\mbox{}}
\fi
\ifx\subparagraph\undefined\else
\let\oldsubparagraph\subparagraph
\renewcommand{\subparagraph}[1]{\oldsubparagraph{#1}\mbox{}}
\fi

% set default figure placement to htbp
\makeatletter
\def\fps@figure{htbp}
\makeatother

\usepackage{etoolbox}
\makeatletter
\providecommand{\subtitle}[1]{% add subtitle to \maketitle
  \apptocmd{\@title}{\par {\large #1 \par}}{}{}
}
\makeatother
\usepackage{booktabs}
\usepackage{amsthm}
\makeatletter
\def\thm@space@setup{%
  \thm@preskip=8pt plus 2pt minus 4pt
  \thm@postskip=\thm@preskip
}
\makeatother
% https://github.com/rstudio/rmarkdown/issues/337
\let\rmarkdownfootnote\footnote%
\def\footnote{\protect\rmarkdownfootnote}

% https://github.com/rstudio/rmarkdown/pull/252
\usepackage{titling}
\setlength{\droptitle}{-2em}

\pretitle{\vspace{\droptitle}\centering\huge}
\posttitle{\par}

\preauthor{\centering\large\emph}
\postauthor{\par}

\predate{\centering\large\emph}
\postdate{\par}
\usepackage[]{natbib}
\bibliographystyle{apalike}

\title{A Student's History of Education}
\author{Frank Pierrepont Graves}
\date{1916}

\begin{document}
\maketitle

{
\setcounter{tocdepth}{1}
\tableofcontents
}
\hypertarget{a-warhorn-historicals-book}{%
\chapter*{\texorpdfstring{\href{https://historicals.warhornmedia.com/}{A Warhorn Historicals Book}}{A Warhorn Historicals Book}}\label{a-warhorn-historicals-book}}
\addcontentsline{toc}{chapter}{\href{https://historicals.warhornmedia.com/}{A Warhorn Historicals Book}}

\includegraphics[width=0.5\textwidth,height=\textheight]{images/warhornlogo.png}

Republished by Warhorn Historicals:
Making historical books available for free online in high quality, readable formats.

We hope this book is a blessing to you. If it is, please \href{https://warhornmedia.com/give}{make a one-time or recurring contribution} right now, sponsor a book from our upcoming list, or volunteer your proofreading or technical skills to help produce more content. Contact us \href{mailto:lucas@beggarsborn.com}{Lucas Weeks} to get involved.

God bless,

---The Warhorn Team

\hypertarget{historical-cover}{%
\subsubsection*{Historical Cover}\label{historical-cover}}
\addcontentsline{toc}{subsubsection}{Historical Cover}

\hypertarget{introduction}{%
\chapter*{Introduction}\label{introduction}}
\addcontentsline{toc}{chapter}{Introduction}

This is the introduction

\hypertarget{the-earliest-education}{%
\chapter{1. The Earliest Education}\label{the-earliest-education}}

\textbf{Outline}

\begin{itemize}
\tightlist
\item
  Even a brief survey of the history of education may greatly broaden one's view.
\item
  Starting with primitive man, we find that his training aims only at the necessities of life, and is acquired informally through the elders and the medicine-men.
\item
  In Oriental education, the next stage in progress, illustrated by India, a traditional knowledge is acquired through \emph{memoriter} and imitative methods.
\item
  While Oriental, Jewish education afforded greater development of individuality, but it was late in organizing schools, \emph{memoriter} in methods, and restricted in content.
\item
  Thus all education before the day of the Greeks was largely \emph{nonprogressive.}
\end{itemize}

\textbf{The Value of the History of Education.} - The History of Education from the earliest times should contribute largely to one's breadth of view and prove a study of the greatest liberal culture. A record of typical instances of the moral, aesthetic, and intellectual development of man in all lands and at all periods should certainly enlarge one's vision and enable him to appreciate more fully the part that education has played in the progress of civilization. Such cultural values may be found even in a limited survey of the world's educational development.

\textbf{It's Treatment in This Book.} - And this is all that will be undertaken here. For, while valuable as a liberal study, the History of Education finds its justification chiefly in the degree to which it functions in the professional training of a teacher, and it will be necessary in a brief treatise to omit or pass over hastily much that might be of interest and value in a more complete account of the development of civilization. Therefore, the amount of space and the perspective afforded the various peoples, epochs, and leaders must here be determined in large measure by the part they have played in the evolution of educational institutions and practices, and by the light their history sheds upon the aim, organization, content, and method of education to-day. At times, too, the history of a single epoch, state, or educational leader will be selected as a type, to the exclusion of others equally important, and treated with considerable intensiveness, instead of describing all sides of the subject with encyclopaedic monotony. Now the first historical epoch to leave a real impress upon modern practice is that of Athens at its height. Hence a mere statement of the salient features of education preceding that period is all that can be afforded in this brief survey. A detailed account of the educational processes used by savage tribes, Oriental nations, and even Judaea may prove interesting and important in other connections, but it must here be largely curtailed.

\textbf{Primitive Education.} There is little to be noted in the training of the young among primitive peoples, save that it is intended largely for the satisfaction of immediate wants - food, clothing, and shelter. Naturally no such actual institution as a school has yet been evolved, but the training is transmitted informally by the parents. The method used is simply that of example and imitation, or, more specifically, `trial and success.' But a more conscious and formal education is given at puberty through the `initiatory ceremonies' (Fig. i). In these rites the youths are definitely instructed by the older men about their relation to the spirits and the totem animals, subordination to the elders, the relations of the sexes, the sacredness of the clansman's obligations, and other traditional usages. Strict silence is enjoined upon them concerning this information, and to impress it upon their minds, and test their endurance, they are required to fast for several days and are often tortured and mutilated. As the savage does not clearly distinguish between himself and the tribe to which he belongs, there is practically no development of individuality, and since the race has not yet learned to treasure its experience in writing, he has no record of past experience and is virtually tied to the present.

\textbf{Oriental Education.} - The nations of the ancient Orient - Egypt, Babylonia, Assyria, China, India, and Persia - may be said to represent the next higher stage in civilization. Their systems of education prepare mostly for vocations, and are not sufficiently advanced to undertake a training for manhood or citizenship. But since a division of labor has now been evolved, the training has become more clearly differentiated and fits for specific occupations. In this way, class divisions, or even castes, have generally arisen in society, and the young people are educated according to the position in life they desire, or are required to fill. As an illustration of this stage of development, we may consider somewhat in detail the social environment and education of India.

\textbf{India: Its Religion and Castes} - In India, largely as a result of the debilitating climate, there was formulated about 1200 B. C. a dreamy philosophy, according to which nothing except Brahma, the one universal spirit, really exists. While men would seem to be temporarily allowed a separate existence of their own, it was held that they should remain inactive as far as possible and seek an ultimate absorption into the great Eternal Spirit. Although somewhat modified by the infusion of Buddhism, between 500 B. C. and 500 A. D., and by the British occupation of the peninsula during the nineteenth century, this mystic and static religion still dominates in India. Connected with it is the caste system, by which the people are divided into four hereditary classes. These are (1) the \emph{brahmins,} or sacerdotal class, which includes all those trained for law, medicine, teaching, and other professional occupations; (2) the warriors, or military and administrative caste; (3) the industrial group; and (4) the \emph{sudras,} or menial caste. Altogether outside the social order are the \emph{pariahs,} or outcasts. The caste system is exceedingly strict. One may fall into a lower caste, but he cannot rise, and loss of caste by one person in a family will degrade all the rest.

\textbf{The Hindu Education.} - Hence Hindu education has always endeavored to fill the pupils with the tenets of their religion, and so prepare them for absorption into the Infinite, rather than for activities in this life, and to preserve the caste system and keep all within the sphere of their occupation. The three upper castes are, therefore, supposed to gain a knowledge of certain sacred works, especially the four \emph{Vedas} or books of `knowledge,' the six \emph{Angas} on philosophical and scientific subjects, and the \emph{Code of Manu,} which is a collection of traditional customs; but few, outside the brahmin class, are ever allowed to take advantage of this opportunity. The warriors are expected to pay more attention to martial exercises, and the industrial caste to acquire through apprenticeship the arts necessary for its hereditary occupations. Sudras, pariahs, and women are generally allowed no education. Except the sudras, all the castes obtain elementary education from a study of the laws, traditions, and customs of the country through the medium of the family, and more recently through village schools held in the open air (Fig. 2). The higher education is largely carried on in brahminic colleges, called \emph{parishads,} and, as also in the case of the elementary work, the teachers have to be brahmins. Since all learning has been preserved by tradition, the chief methods of instruction are those of memorizing and imitation. Even the later texts are so written as to be easily committed, and the lines are sung aloud by the pupils until they have memorized them. Writing is learned by imitating the teacher's copy on the sand with a stick, then on palm leaves with a stylus (Fig. 2), and finally on plane leaves with ink.

\textbf{Effect of the Hindu Education.} - Hence, among the Hindus education is forbidden to ninety-five per cent of the population, and, as far as it does exist, it is a mere stuffing of the memory. It concerns itself but little with mental culture or with preparation for real living. The brahmins have handed down considerable traditional learning, grammar, phonetics, rhetoric, logic, `Arabic' notation, algebra, astronomy, and medicine, but new knowledge of any sort is barred. The Hindus still plow with sticks of wood, and their crops are harvested and threshed by devices equally primitive. They bake bricks, work metals, and weave cloth, but with the same kind of appliances that were used by their remote ancestors. Until recently, they have been greatly lacking in ambition, self-reliance, and personal responsibility, and have not yet come to any feeling of solidarity or national unity. To them prosperity and progress are foreign ideas.

\textbf{India as Typical of the Orient.} - The other countries of the ancient Orient never fixed their social classes in so hard and fast a manner, and have never included so elaborate a philosophy among the products of their culture. But India may well be considered broadly typical of the stage of development in the Orient. Certain common features appear in the education of all the nations there. In the system of each, the classes below the sacerdotal or priestly are given little intellectual education, and the women none at all, but both are trained by apprenticeship in their vocations. Actual schools, both elementary and higher, have been instituted; and the latter, except in China, are conducted at temples or priestly colleges by members of the sacerdotal class. The educational content is naturally traditional. It is, for the most part, ensured against change by being embalmed in sacred books, such as the \emph{Vedas.} The educational method consists largely in the memorizing of the text and imitation of the copy set, and little attempt is made to give a reason for the customs and traditional knowledge taught. Hence, while individuality has begun to emerge, it is suppressed by every agency possible; and, although these peoples have largely overcome the primitive enslavement to nature and the present, they are completely in bondage to the past.

\textbf{Jewish Education.} - The Jews are classed among the nations of the Orient, but they formulated loftier aims and have exerted more influence upon modern ideals in education. While their theology greatly developed in the course of their history, from the first they held to an ethical conception of God, and the chief goal of their education was the building of moral and religious character. Not until after the Babylonish captivity (586-536 B. C), however, did they establish actual schools. Before that, children were given an informal training in the traditions and observances of their religion by their parents. But they brought back from Babylon the idea of institutions for higher training and started such schools through their synagogues. In the second century B. C. the founding of elementary schools also began, and eventually the Jews made education well-nigh universal. The beneficial effect of this training is seen in the respect shown by the Jews for their women, their kind treatment of children, and their reverence for parents. The defects of their education appear in the stereotyped and formal way in which the religious material came to be interpreted, and the consequent hostility to science and art, except as they threw light on some religious festival or custom. Although appeal was made to various types of memory, systems of mnemonics devised, and other good pedagogical features suggested, their methods of instruction were largely \emph{memoriter.} The Jewish system of education, as a whole, afforded a greater development of personality than that of the other Oriental nations, and through it have been spread some of the world's most exalted religious conceptions. Nevertheless, it did not depart much from its traditions and the past, and to this extent it may be classed with the training of the primitive tribes and of the Oriental nations as predominantly \emph{non-progressive.}

Supplementary Reading

For general works, see Graves, F. P., \emph{History of Education before the Middle Ages} (Macmillan, 1909), chaps. I-XI; Monroe, P., \emph{Text-book in the History of Education} (Macmillan, 1905), chaps. I-II. A general interpretation of the evolution of education in savagery and barbarism is also given in Laurie, S. S., \emph{Pre-Christian Education} (Longmans, Green, 1909), pp.~1-207; Morgan, L. H., \emph{Ancient Society} (Holt, 1907), Part I; and Taylor, H. O., \emph{Ancient Ideals} (Macmillan, 1913), vol.~I, chaps. I-V. An illustration of primitive training of especial interest to American students is found in Spencer, F. C., \emph{Education of the Pueblo Child} (Columbia University, Department of Philosophy and Psychology, vol.~7, no. 1); and a detailed description of the puberty rites of a variety of savage tribes, in Webster, H., \emph{Primitive Secret Societies,} (Macmillan, 1908), chaps. I-V. A more complete account of the Hindu philosophy and education appears in Dutt, R. C, \emph{Civilization of India} (Dent, London), and Taylor, H. O., \emph{Ancient Ideals} (Macmillan, 1913), vol.~I, chaps. Ill and IV. A systematic statement of the Jewish training has been adapted from a German work, in Leipziger, H. M., \emph{Education of the Jews} (New York Teachers College, 1890), and a more detailed account worked out in Spiers, B., \emph{School System of the Talmud} (Stock, London, 1898).

\hypertarget{the-education-of-the-greeks}{%
\chapter{2. The Education of the Greeks}\label{the-education-of-the-greeks}}

\textbf{Outline}

\begin{itemize}
\tightlist
\item
  The Spartan training was intended to serve the state by making warriors, and little attention was paid to intellectual education.
\item
  At first the Athenian education was also mainly concerned in serving the state. For the earliest stage of the boy's education, there were schools of two types, - one for intellectual training, as well as one for physical; from fifteen to eighteen a more advanced physical training was given; and then, for two years, a preparation for military life.
\item
  After the Persian wars, the Athenians adopted ideals of education affording a larger recognition of individualism. The sophists introduced the new educational practices, and went to an extreme in their individualism.
\item
  The systematic philosophers, - Socrates, Plato, and Aristotle, tried to mediate the outworn institutional education and the extreme individualism. Socrates held that the sophistic `knowledge' was only `opinion,' and that the more universal knowledge could be reached in every person by stripping off his individualistic opinion.
\item
  But Plato maintained that only the intellectual class could attain to knowledge. For them he formulated a new course of study, in addition to that in vogue, consisting of mathematical subjects and dialectic. Aristotle held that the training for every one before seven should be bodily; up to fourteen, the irrational soul should be trained; and until twenty-one, the rational. While Plato and Aristotle had little effect upon educational practice at the time, they have since greatly influenced education.
\item
  After Aristotle, there arose individualistic schools of philosophy and formal schools of rhetoric, and out of them universities sprang up. Then Greek culture and education spread throughout the world.
\end{itemize}

\textbf{Progressive Nature of Greek Education.} - Real educational progress began with the Greeks. In their training gradually appeared considerable regard for individuality. They were the first people whose outlook seems to have been toward the future rather than the past, and they first made a serious attempt to promote human development in accordance with a remote ideal progressively revealed. As a result, they not only gave a wonderful impetus to educational practice in their own time, but ever since then the world has had constant recourse to them for inspiration and counsel. While this intellectual emancipation did not appear to any extent before its development among the Athenians in the middle of the fifth century B. C, well-planned systems of education existed in Greece several centuries before this and paved the way for the system in Athens during the Age of Pericles.

\textbf{Spartan Education: Its Aim and Early Stages.} - Among the states of ancient Greece, Sparta possessed the earliest education of which we have any extended information. Its citizens dwelt in the midst of hostile peoples they had subjugated, and this made it necessary to produce a race of hardy and patriotic warriors. Strength, courage, and obedience to the laws were held as the aim of education. The Spartan educational system was intended to serve the state, and the rights of the individual were given little or no consideration. State control began with birth. The infant was immediately inspected by a council of elders, and, if he were sickly or deformed, he was exposed to die in the mountains; but if he appeared physically promising, he was formally adopted by the state and left with his mother for rearing until seven. At that age the boys were placed in charge of a state officer and ate and slept in a kind of public barracks. Here their life became one of constant drill and discipline. In addition to hard beds, scanty clothing, and little food, they were given a graded course in gymnastics. Besides ball-playing, dancing, and the \emph{pentathlum} - running, jumping, throwing the discus, casting the javelin, and wrestling - the exercises included boxing, and even the brutal \emph{pancratium,} in which any means of overcoming one's antagonist - kicking, gouging, and biting, as well as wrestling and boxing - was permitted.

The Spartan boys, however, received only a little informal training in the way of intellectual education. They simply committed to memory and chanted the laws of Lycurgus and selections from Homer, and they listened to the conversation of the older men during the meals at the common table, and were themselves exercised in giving concise and sensible answers to questions put to test their wisdom. Every adult was also required to choose as his constant companion or `hearer' a youth to whom he might become an'inspirer.'

\textbf{Training in Youth and Manhood: Results.} - When a youth reached eighteen, he began the distinctive study of warfare. For two years he was trained in the use of arms and skirmishing, and every ten days had his courage and his physique tested by being whipped before the altar of Artemis. Then he regularly entered the army, and for ten years guarded some border fortress and lived upon the coarsest of fare. When he became thirty, he was considered a man and forced to marry at once, but even then he could visit his wife only clandestinely and was still obliged to live in common with the boys and assist in their training.

The education of women was very like that of the men. While the girls were allowed to live at home, they were given a similar physical training in the hope that they would become the mothers of sturdy sons. Thus the Spartan education was shaped entirely with reference to the welfare of the state. Their educational system served well its purpose of creating strong warriors and devoted citizens, but it failed to make for the highest manhood. Sparta developed practically no art, literature, or philosophy, and produced little that tended to promote civilization. She has left to the world little but examples of heroism and foolhardiness alike.

\textbf{Old Athenian Education: Its Aim and Early Training.} - For many centuries the Athenian education was not unlike the Spartan in promoting the welfare of the state without much consideration of individual interests. But even in early days Athens felt that the state was best served when the individual secured the most complete personal development. Hence, the Athenian boys began to receive at seven years of age two kinds of training - (1) the \emph{pentathlum} and other physical exercises in the \emph{palaestra} (Fig.3) or exercising ground, and (2) singing and playing upon the flute or lyre, and reading and writing at the \emph{didascaleum} (Fig. 4.) or music school. After the boy had learned his letters by tracing them in the sand, he was taught to copy verses and selections from well-known authors, at first upon wax-tablets with a stylus, and later upon parchment with pen and ink. It was, moreover, necessary for the pupils in singing to be taught the rhythm and melody, and to understand the poem so as to bring out its meaning. Hence the explanations and interpretations given by the teachers brought in all the learning of the times, and the moral and intellectual value of the studies must have been much greater than would be suggested by the meagerness of the course. Some moral training and discipline were also given the boy by a slave called the \emph{paedagogus,} who conducted him to school and carried his lyre and other appurtenances. This functionary was often advanced in years or incapacitated for other duties by physical disability.

\textbf{Training for the Youth.} - At fifteen the Athenian boy might take physical training of a more advanced character at one of the exercising grounds just outside Athens, which were known as \emph{gymnasia.} He was now ephebic course permitted to go wherever he wished and become acquainted with public life through first-hand contact. When eighteen the youth took the oath of loyalty to Athens, and for two years as an \emph{ephebus} or cadet continued his education with a course in military duties. The first year he spent in the neighborhood of Athens and formed part of the city garrison, but in the second year he was transferred to some fortress on the frontier. At twenty the young man became a citizen, but even then his training continued through the drama, architecture, sculpture, and art that were all about him.

\textbf{Effect of the Old Athenian Education.} - Little attention was, however, given by the Athenians to the education of woman. It was felt that her duties demanded no knowledge beyond ordinary skill in household affairs. With this exception, the Athenian education was superior to the Spartan in allowing greater opportunity for individual development and in furnishing a more rounded training. Nevertheless, until about the middle of the fifth century B.C., while differing considerably in degree from Sparta, Athens may be grouped with that country as adhering to the `old' education, where the individual was subordinated to the good of the social whole.

\textbf{Causes and Character of the New Athenian Education.} - This characterization is, of course, in contrast to Greek education in the `new' period, which is represented by Athens alone. This later type of education was probably somewhat the result of the gradual rise of democratic ideals in Athens, but a more immediate set of factors grew out of the Persian wars (492-479 B. C). This extended conflict with a powerful Oriental people, possessing a well-organized but widely different body of traditions tended to broaden the views of the Athenians greatly, and the ensuing political and commercial intercourse with a variety of dependent states and nations in the Delian League, together with social contact with the foreigners from every land that were thronging the streets of Athens, led even more directly to a reconstruction of practices and beliefs. A rapid transition in the old traditions took place and society seems for a time to have been sadly disorganized. The old was shattered, and while new ideals were being constructed, a groping ensued. Although the latitude given the individual was destined, as always, to produce progress in the long run, and was of great ultimate service to the world, more immediately a low ebb in morals at Athens resulted. Individualism ran riot. Education reflected the conditions of the period. Its ideals became more and more individualistic. The times demanded a training that would promote the happiness of the individual with little consideration for the welfare of the state as a whole. The old education seemed narrow and barren of content; and there arose a desire for all sorts of knowledge that might contribute to one's advancement, whether it increased his social usefulness or not. Skill in debate and public speaking was especially sought, because of the unusual opportunity for personal achievement in politics.

\textbf{The Sophists and Their Training.} - To meet these new demands, a set of teachers known as the \emph{sophists} came into prominence. They professed to train young men for a political career, and some of them even claimed to teach any subject whatsoever, or how to defend either side of an argument. These pretensions, together with their charging a fee for their services, contrary to Athenian custom, seriously offended the more conservative of the citizens of Athens. But many of the first sophists afforded an honest and careful training. The effect of their teaching was especially felt by the adolescents in the \emph{gymnasium} stage of education, since they were ambitious to distinguish themselves politically. The physical training that had hitherto dominated the gymnasium course gave way to a study of grammatical and rhetorical subtleties, and whenever a sophist appeared in the street, market-place, or house, the young men crowded about him to borrow from his store of experience and wisdom, and acquire his method of argument. To a less degree the same influence was felt in the lower schools and by the cadets and younger citizens. The exercises of the palaestra were no longer as rigorous, and existed for the sake of individual health and pleasure rather than for the making of citizens. The literary work of the didascaleum came to include, besides the Homeric epics, a wide range of didactic, reflective, and lyric poetry, with a superabundance of discussions. In music the old patriotic and religious songs sung to the simple Doric airs and accompanied upon the sevenstringed lyre, were replaced by rhythms of great difficulty, like the Lydian and Phrygian, and by complicated instruments of all sorts.

\textbf{Their Extreme Individualism.} - All this inroad upon the time honored curriculum shows how fully the sophists embodied the individualism of the times. Although they held no body of doctrine common to them all, they were generally at one in their position of extreme individualism. They often went so far as to insist that there could not safely be any universal criteria in knowledge or morals; that no satisfactory interpretation of life could be made for all, but that every fact and situation should be subject to the judgment of the individual. No doubt the formula attributed to Protagoras, ``Man (i. e. the individual) is the measure of all things, both of the seen and the unseen,'' would have expressed the attitude common to most of them. They but carried to its legitimate conclusion the complete reaction from the old ideal of subordination of the individual to the state.

\textbf{The Reactionaries and the Mediators.} - Meanwhile, the conservative element was making its usual attempt to adjust the unsettled conditions by suggesting a return to the old. Various schemes had been advanced, even before the sophists had come into prominence. Of these the most complete plan was that of Pythagoras (about 580-500 B.C.). By adopting an analogy from the `harmony' of the celestial bodies and from the relation of the powers in the individual to each other, he arranged a definite hierarchy in society, so that each member should have his proper place, and complete harmony and social order should ensue. As the influence of the sophists began to be felt, later representatives of the reactionary movement, such as the matchless caricaturist, Aristophanes (445-380 B.C.), began to appear and inveigh against the new conditions. But the social process can never move backward, and reconstruction on some higher plane was needed to overcome the destructive tendencies of the times. To furnish this, was the task set themselves by Socrates, Plato, and Aristotle. Like the sophists, they recognized that the traditional beliefs and sanctions, the old social order, and the former ideals and content of education, had been outlived, and that the individual could not find truth and morality through an institutional system. At the same time they felt that the extreme individualism of the sophists was too negative a basis upon which to build, and that a more socialized standard of knowledge and morality must be sought.

\textbf{The Method of Socrates.} - This mediating effort was begun by Socrates (469-399 B. C). While he started with the formula of Protagoras, he maintained that the `man' indicated thereby was not the individual, but mankind as a whole. It is not the peculiar view of any individual that represents the truth, but the knowledge that is the same for everyone. The former, which the sophists considered `knowledge,' Socrates held to be only `opinion, and declared that the reason men think so differently is because each sees but one side of the truth. He believed that everyone could get at universal knowledge by stripping off individual differences and laying bare the essentials upon which all men are agreed. He conceived it to be the mission of the philosopher or teacher to enable the individual to do this, and he endeavored to deal with the mind of all those with whom he came in contact, so that they would form valid conclusions. By his method, known as the \emph{dialectic,} or 'conversational,' he first encouraged the individual to make a definite statement of his belief, and then, through a set of clever questions, caused the person to develop his thought, until he became so involved in manifest contradictions that he was forced to admit that his view had been imperfectly formed. He thus caused the individual to see that the view he had first expressed was mere `opinion' and but a single phase of the universal truth. As Socrates further held that morality consists in right knowledge and made no distinction between the knowledge of an action and the impulse to perform it, he strove through his methods of developing knowledge to harmonize the individual welfare with that of the social group.

\textbf{Plato's System of Education for the Three Classes of Society.} - But the believers in the old traditions and institutional morality felt that Socrates was atheistic and immoral. They persuaded Athens to give him the hemlock, and thus destroyed the man who might have proved her savior. A pupil, Plato (427-347 B. C), undertook to continue his work, but his aristocratic birth and temperament caused him to underestimate the intelligence of the masses. He held that they were incapable of attaining to `knowledge' - that they possessed only `opinion.' In his most famous dialogue, \emph{The Republic,} he endeavors to show that the ideal state can exist only when the entire control of the government is entrusted to the `philosophers,' or intellectual class, who alone possess `real knowledge.' Those who are to compose the three classes of society Plato would have selected during the educational process on the basis of their ability. For all boys up to eighteen years of age he prescribes an education similar to that in vogue in the palaestra, didascaleum, and gymnasium, except that he would somewhat expurgate the literary element, and would confine the musical training to the simpler melodies and instruments. The youths who prove capable of going beyond this lower education are next to take up the cadet training between eighteen and twenty, but those who are incapable of further education are to be relegated to the industrial class. During the cadet period are to be determined those capable of going on with the higher education of philosophers, while those who here reach their limit become members of the military class.

As Athenian education did not extend beyond the twentieth year, Plato is here obliged to invent a new course of study that will enable the future philosophers to acquire the habit of speculation. This additional course, he declares, should also be graded, in order that a further test of intellectual and moral qualities may be made. Arithmetic, plane and solid geometry, music, and astronomy, are to occupy the first ten years of the course. These subjects, however, are not to be studied for calculation or practical purposes of any sort, but entirely from the standpoint of theory or the universal relations underlying them, since only thus can they furnish a capacity for abstract thought. After this, at thirty, the young men who can go no further, are to be placed in the minor offices of the state, while those who have shown themselves capable of the study of dialectic, go on with that subject for five years longer. It then becomes the duty of these highest philosophers to guide and control the state until they have reached the age of fifty, when they may be allowed to retire.

\textbf{The Weakness of Plato's System.} - Thus, where Socrates found the basis of universal truth in everyone, Plato held that only one class of people, the most intellectual, could attain to real knowledge. He, therefore, maintained that the philosophers should absolutely guide the conduct of the state, and that education should be organized with that in view. Plato's ideal state would thus become a sort of intellectual oligarchy, and in a way was a return to the old principle of subordinating the individual to society. \emph{The Republic} thus quite neglected human will as a factor in society and assumed that men can be moved about in life like pieces upon the chess board. Plato failed to see, too, that each individual really possesses all human characteristics. The workers have reason, and the philosophers have passions, and a human being is not a man unless all these functions are his. But even if his scheme had been a happy one, the treatise provided no method of evolution from current conditions, and if it were further granted that this order of things could be established at once. Plato put the ban upon all innovation or change, and so closed the door to progress.

Hence \emph{The Republic} was viewed as a visionary conception, and had no immediate effect upon education or any other institution of Athens. So in his declining years, without denying \emph{The Republic} as ideal, he wrote the more practical dialogue known as \emph{The Laws.} In it he welded elements from the educational systems of Sparta and older Athens, and reverted to traditions and ideals not dissimilar to the doctrines of Pythagoras. He replaced the philosophers with priests, an hereditary ruler, a superintendent of education, and various other officials; and the course of study reached its height with the subject of mathematics, while dialectic was not mentioned.

\textbf{His Influence upon Educational Theory and Practice.} - Thus the efforts of Socrates, as continued by Plato, to obtain the benefit of the growing individualism for society and education without disrupting them, had seemingly come to naught. Nevertheless, Plato has had considerable influence upon the thought and practice of men since the Greek period. The ideal society where everything is well managed and everyone is in the position for which nature intended him, has ever since the day of \emph{The Republic} been a favorite theme for writers, as witness More's \emph{Utopia} and the \emph{New Atlantis} of Bacon. A specific movement that shows the impress of Plato, as we shall see later, is the formulation of the more advanced studies of the mediaeval `seven liberal arts' under the name of the `quadrivium.' It is even possible that the whole conception of `liberal' studies, and so the doctrine of `formal discipline' (see p.~182), may be traced back to Plato's idea that the mathematical subjects in the course for philosophers should never be studied from a practical point of view. On the whole, Plato has been a factor in educational theory and practice that cannot be overlooked.

\textbf{Aristotle's Ideal State and Education.} - A more practical attempt to unify the new with the old in Athenian society and education was made by Aristotle (386322 B. C), the pupil of Plato. From his father, the court physician at Macedon, and from his study under Plato, Aristotle obtained an excellent scientific training, which is evident in the way he approaches his problems. It is in his \emph{Politics} especially that he discusses the ideal state and the training of a citizen. His method of investigation to determine the nature of this ideal state is inductive, and before formulating his conception of it, he makes a critical analysis of Plato's \emph{Republic} and \emph{Laws,} and analyzes the organization of many other states, both ideal and actual. He concludes that a monarchy is theoretically the best type of government, but that the form most likely to be exercised for the good of the governed is the democracy. He then considers in detail the best natural and social conditions for a state. Among these practical considerations is the proper education to make its citizens virtuous.

Since virtue is of two kinds, moral or practical, and intellectual or speculative, and the former is merely the stepping-stone to the latter, the education needed for the virtue of the state must not, like that of Sparta, be purely a training for war and practical affairs. In marking off the periods of education, Aristotle holds that ``the care of the body ought to precede that of the soul, and the training of the impulsive side of the soul ought to come next; nevertheless, the care of it must be for the sake of the reason, and the care of the body for the sake of the soul.'' The development of the body he wishes to start even before birth by having the legislator ``consider at what age his citizens should marry and who are fit to marry.'' Also he deems it necessary to sanction the usage of his time of `exposing' (see p.~13) all deformed and weakly children. However, his advice concerning the food, clothing, and exercise of children is humane and in keeping with the best modern hygiene.

The training of the body is a preparation for the formal schooling, which is to last from seven to twenty-one. This is divided into two periods by puberty, the first to be devoted to the training of the impulsive or irrational side of the soul, and the second to that of the rational side. Education, he claims, should be public, as in Sparta, for it is the business of the state to see that its citizens are all rendered virtuous. However, the industrial classes, not being citizens, have no need of education, and women are to be limited in the scope of their training. The course of study for the irrational period is largely the same as that in use at Athens, - gymnastics, music, and literary subjects, although he recommends some reforms. Gymnastics is intended for self-control and beauty of form, and the making of neither athletes nor warriors should be the object, since the training of the former exhausts the constitution, and that of the latter is brutalizing. The literary subjects, which with Aristotle includes drawing, as well as reading and writing, are not to be taught merely for utilitarian reasons. Music is to be used not so much for relaxation or intellectual enjoyment as for higher development. Since melodies that afford pleasure are connected with noble ideas, and those which give us pain are joined to debased ideas, the study of music ``cultivates the habit of forming right judgments, and of taking delight in good dispositions and noble actions.'' Another moral effect of music is that it produces \emph{katharsis} or `purification'; that is, by arousing in us pity and fear for humanity at large, it lifts us out of ourselves and affords a safe vent for our emotions.

Such was to be the training for the body and for the irrational period, but how Aristotle would have advised that the education of the rational soul be carried on can only be surmised, since the treatise breaks off suddenly at this point. It is probable that it would have included a higher training in mathematical subjects and dialectic similar to that advocated by Plato, and, from Aristotle's own predilections, he would have been likely also to add some of the physical and biological sciences.

\textbf{The Permanent Value of His Work.} - Thus Aristotle, like Plato, endeavored to work out the harmonizing of individual with social interests by the creation of an ideal state, and he similarly failed to answer the demand of the times. His work was much less visionary than \emph{The Republic,} but he did not fully recognize that the day of the small isolated states of Greece, with their narrow prescriptions for patriotism and social order, had passed forever. Hence he hoped to achieve some reform by departing but little from existing conditions and reading a philosophy into them, and this bondage to the times prevented his educational system from making any advance beyond that of Plato. But while Aristotle had little effect upon the society of the times, his works have since been considered of great value, and the methods that he formulated have been most important. He not only started, or made the first great contributions to a number of sciences, but he crystallized the laws of thought itself. Also, as instruments to assist in fashioning the various sciences, Aristotle invented a complete system of terminology, and created such pairs as `matter' and `form,' `mean' and `extreme,' and `cause' and `effect,' and such convenient expressions as `principle,' `maxim,' `habit,' and `faculty.' A more important effect of Aristotle's ideas has been that upon the formulation of doctrine in the Christian Church. After the spread of Mohammedanism, which had largely absorbed the Aristotelian principles, the Church, though at first bitterly opposing them, finally found it impossible to suppress them, and began to clothe her own doctrine in their dress. The greatest of the scholastics began to study Aristotelianism, and soon made it the effective weapon of the Church by reducing all human knowledge to a finished Aristotelian system with theology at the top.

\textbf{The Post-Aristotelian Schools of Philosophy.} - But the harmonizing attempt of Aristotle was fruitless. Like Socrates and Plato, he failed to reconcile with the old and settled order the ever-expanding movement toward individualism. Thus all efforts to control the individualistic and disintegrating tendencies of the times were in vain, and the conquest of the Greek states by Philip of Macedon (358-338 B. C) was only symptomatic of the complete collapse of corporate life and the inability to reconstruct it successfully. All possibility of social unity disappeared, and philosophy no longer considered the individual from the standpoint of membership in society. It was occupied no further with the harmonization of the individual and the state, but concerned itself with the welfare of the individual and the art of living. Individualism was completely triumphant, and education was considered simply as a means to personal development or happiness, without regard to one's fellows. The new theories of life and education were formulated by such schools of philosophy as the Epicureans, Stoics, and Skeptics, which kept themselves far removed from society. None of these `schools' could be so termed in the sense of offering an education, but rather in the modern usage of a group of adherents to certain teachings. They spent their energy, for the most part, in interpreting, elaborating, and lauding the original teachings of the founders, and with them a stereotyped dogmatism took the place of philosophy.

\textbf{The Schools of Rhetoric.} - But these schools were not the only outcome of the teaching of the sophists. Just as they came about gradually from the speculative tendencies of the sophists as developed through certain famous philosophers, there likewise grew up more directly from the sophistic efforts to train young men in rhetoric and public speaking a multitude of rhetorical schools. In these a formal study was made of oratory and the knowledge of the day. Their professed object was to make successful men of the world, and, although they at first included such reputable and influential schools as that of Isocrates (436-338 B. C), they laid little claim to teaching anything solid or profound, much less to forming any philosophic habits. They succeeded in spreading a popular education among a people that had lost all hope of a political life, but they soon degenerated into the use of narrow and formal methods. The later rhetoricians attempted to hasten oratorical training and preparation for life, by teaching their pupils ready-made speeches and dialogues, together with a general knowledge of current questions. Nevertheless, these schools flourished for several centuries and closely rivalled those of the philosophers.

\textbf{The Hellenic Universities.} - From these two classes of schools, the philosophical and the rhetorical, the fame of Athens spread rapidly, and from the fourth century B. C. onward the number of young men from all over the civilized world who came there to study steadily increased. Before the close of the century the old cadet training of Athens was united with this intellectual education, and there sprang up a regular institution or university, which the young Athenians and students from outside might attend. Before long, the Hellenic world boasted other universities, such as those at Rhodes, Pergamon, Alexandria, and Rome. Until almost 300 A. D. Athens remained the chief intellectual center of civilization, and attracted students from all parts of the Roman Empire. Gradually, however, the higher education there tended toward the study of rhetoric alone and artificiality grew apace. In consequence, Alexandria came to displace Athens as the center of culture, and her university became the leading one of the world. Here the various philosophic and religious sects gathered to study and discuss, and the abstract Greek philosophy united with the more concrete beliefs of the Orient, especially Zoroastrianism, Judaism, and Christianity. Thus there flourished here the various systems of religious philosophy known collectively as `Hellenistic,' such as Neopythagoreanism, Neomazdeism, Philonism, Gnosticism, and Neoplatonism. Considerably before this, too, there had developed at Alexandria the Ptolemaic theory of the universe. Other noted investigations, like those of Euclid in geometry, Archimedes in physics, Eratosthenes in astronomy, and Diophantus in algebra, also bore witness to the intellectual activity of this university.

\textbf{Extension of Hellenic Culture.} - It can thus be seen that the political downfall of Athens had only prepared the way for a larger intellectual influence. As Alexander extended his yoke over one Eastern country after another, he had carried with him all the culture of Greece, and within a century of his death the whole Orient was dotted with Greek gymnasia, stadia, and theaters, and saturated with Greek literature, art, philosophy, and education. Similarly Rome, which had come somewhat into contact with Greece before conquering her, had been tinctured with Greek life and learning; and, after her absorption of Macedon and Greece, she fell under the spiritual thrall of the subjugated people. The history of Greek civilization and education was so intermingled with the Roman that it can scarcely be distinguished from it. The Greek schools of philosophy and rhetoric were continued in Rome, Roman youths made up a great body of the attendance at the universities of Athens and Alexandria, and the Roman emperors did much for the support and extension of the work in these institutions. Hence from the Greeks have developed some of the most advanced intellectual and aesthetic ideas that civilization has known.

\textbf{Supplementary Reading}

Graves, \emph{Before the Middle Ages} (Macmillan, 1909), chap. XII; Monroe, \emph{Text-book} (Macmillan, 1905), chap. III. See also Laurie, \emph{Pre-Christian Education} (Longmans, Green, 1900), pp.~208-318. Davidson, T., in his \emph{Aristotle} (Scribner, 1896), develops the periods of Greek education in chronological order, and his \emph{Education of the Greek People} (Appleton, 1903) gives the social setting of its development. A most scholarly and brilliant work is Freeman, K. J., \emph{Schools of Hellas} (Macmillan, 1907), which is illustrated by vasescenes and other reproductions of Greek education. Bosanquet, B., \emph{The Education of the Young in Plato's Republic} (Cambridge University Press, 1908), Nettleship, R. L., \emph{Theory of Greek Education in Plato's Republic} (See Evelyn Abbott's \emph{Hellenica,} Longmans, Green, 1908), and Burnet, J., \emph{Aristotle on Education} (Cambridge University Press) afford a good interpretation of the theorists mentioned; while Capes, W. W., in the \emph{University Life in Ancient Athens} (Harper, 1877), and Walden, J. W., in the \emph{Universities of Ancient Greece} (Scribner, 1909), furnish a lively description of the students and professors.

\hypertarget{the-education-of-the-romans}{%
\chapter{3. The Education of the Romans}\label{the-education-of-the-romans}}

\textbf{Outline}

\begin{itemize}
\tightlist
\item
  The contribution of the Romans to progress was largely due to their absorption of Greek culture, but their primitive training had an influence in itself. This was mostly civic and practical, and was given informally in the family and the forum.
\item
  Through amalgamation with the Greek, Roman education maintained three grades of schools: (i) the elementary school or \emph{ludus,} (2) the `grammar' school, and (3) the rhetorical school. Beyond the education of these schools, a young Roman might attend a university.
\item
  Schools were gradually subsidized by the emperors, but education eventually deteriorated into a formal qualification for senatorial rank. The practical Romans, however, created a universal empire and legal system, a universal religion, and other institutions for modern society.
\end{itemize}

\textbf{Roman Education Amalgamated with Greek.} - The name of Rome is still suggestive of power and organization. These characteristics seem to have been innate; but the significance of Roman development to the history of progress and education was largely due to the fact that, in her spread over the civilized world, the Eternal City amalgamated the Greek civilization with her own. Until then her ideals of life, while effective in conquest, had been narrow and little adapted to the development of individuality or of cosmopolitanism. Unconsciously realizing the need of broader ideals, she absorbed those of Greece. But Rome could not be Hellenized without making some contributions to the result from her own genius, and for that reason it is important to learn something of Roman civilization and education, crude as they were, before they came into contact with Greek culture.

\textbf{Early Education in Rome.} - In the early days Rome was animated by intense patriotism and love for military life, and felt that each citizen was bound to merge his identity in that of the state. In the surrender of individuality they were, to be sure, not unlike the Spartans, although they believed that this subordination should be brought about voluntarily rather than by compulsion of law and state organization. But, with such a love as theirs for mere material achievement, the Athenian ideal of a full and harmonious development of one's whole nature could scarcely be expected to make any appeal. They looked not for harmony, proportion, or grace, but for stern utility. They were sedate, grave, and serious, and their education was practical, prosaic, and utilitarian.

Until the Greek institutions began to be adopted, schools did not exist in Rome, except possibly the \emph{ludus} or elementary school. During this pristine period education consisted in a practical training in Roman ideals and everyday living conducted largely through the family. In childhood the boys and girls alike were given a physical and moral training by their mother, but, as the boy grew older, he went more in the company of his father, and learned efficiency in life informally through his example and that of the older men, while the girl was taught at home by her mother. If the boy belonged to a patrician family, he might acquire much knowledge concerning Roman custom and law by hearing his father advise and aid the family \emph{clients,} or `dependents,' and by attending banquets with him. He might also receive an apprenticeship training from his parent or some other older man in the profession of soldier, advocate, or statesman. In case he was born in a less exalted station, he might learn his father's occupation at the farm or shop. The girl, whatever her social status, was trained by her mother in the domestic arts, especially in spinning and weaving wool. Through their parents children probably learned to read and write; and they committed to memory stories of Roman heroes, ballads, martial and religious songs, and the \emph{Twelve Tables} of national laws, after these had been codified (451 B. C). Physical exercise was secured largely by games, which were mostly in imitation of future occupations, and gymnastics were employed simply as training for war. The usages of home and public religion also played an important part in the education of the young Romans, especially since almost every activity in life was presided over by some deity, whom it was necessary to propitiate when engaging in it. Thus education in early Rome was practical, and, to some extent, occupational. It was intended to produce efficiency as fathers, citizens, and soldiers. It consisted in training the youths to be healthy and strong in mind and body, and sedate and simple in their habits; to reverence the gods, their parents, the laws, and institutions; and to be courageous in war, and familiar with the traditional agriculture, or the conduct of some business. It did produce a nation of warriors and loyal citizens, but it inevitably tended to make them calculating, selfish, overbearing, cruel, and rapacious. They never possessed either lofty ideals or enthusiasm. Their training was best adapted to a small state, and became unsatisfactory when they had spread over the entire Italian peninsula. The golden age of valor and stern virtue had then largely departed, and they began unconsciously to seek a more universal culture. While such a people regarded the Greeks as visionary, just as the Greeks looked upon them as barbarians, they felt instinctively that only by absorption of the Hellenic ideals could their cosmopolitan ambitions be carried out. On the other hand, it was through the organization which the Romans were able to furnish, that the great ideals formulated by the Greeks were destined to be rendered effective and to become a matter of value and concern to civilization ever since.

\textbf{The Absorption of Greek Culture.} - There was a gradual infiltration of Greek culture into Rome from very early days. This received a great impulse through the conquests of Alexander (334-323 B. C) and the absorption of Macedon by Rome (168 B.C.), but it was not until about half a century after Greece itself had become a Roman province (146 B. C.), that the Greek educational ideals and institutions can be said to have been completely absorbed by Rome. This new type of education was thus well established early in the first century B.C. It may be said to have remained almost unmodified until toward the end of the second century A. D., when political conditions at Rome became most unstable and the period of degeneracy set in. During these three centuries of Hellenized Roman education, three grades of schools resulted from the amalgamation. They were the (1) \emph{Iudus} or school of the \emph{litterator,} as the lowest school was called; (2) the `grammar' school, taught by a \emph{grammaticus} or \emph{litteratus;} and (3) the schools of rhetoric and oratory, which furnished a somewhat higher education.

\textbf{The Ludus.} - The ludus, or lowest school, may possibly have existed before the process of Hellenization even began, but if it did, it must have been intended simply to supplement the more informal training of the home. Whenever originated, it probably taught at first only reading, writing, and rudimentary calculation, as in the family, through the medium of historical anecdotes, ballads, religious songs, and the \emph{Twelve Tables.} But as the Greek influence crept in more and more, the literary content was somewhat extended. About the middle of the third century B. C, Livius Andronicus translated the \emph{Odyssey} into Latin; and a number of epics, dramas, and epigrams were soon composed after Greek models. These works, in whole or part, were introduced into the curricula of the \emph{ludi,} and by the beginning of the first century B.C., the \emph{Twelve Tables} had been displaced by the Latinized \emph{Odyssey} of Andronicus. The methods of instruction were \emph{memoriter} and imitative. The names and alphabetic order of the letters were first taught without any indication of their significance or even shape, and all possible combinations of syllables were committed before any words were learned. Reading and writing were then taught by dictation, and, in tracing the letters on wax-tablets with the stylus (Fig. 5), the hand of the pupil was at first guided by the teacher. Calculation was learned by counting on the fingers, by means of pebbles, or upon the abacus, and eventually sums were worked upon the tablets.

Methods so devoid of interest were naturally accompanied by severe discipline. The rod, lash, and whip seem to have been in frequent use, and the names ordinarily applied to schoolmasters in Latin literature are suggestive of harshness and brutality. Moreover, a fresco found at Herculaneum depicts a boy held over the shoulders of another, with the master beating the victim upon the bare back (Fig. 6). Under these circumstances, no real qualifications were required of the teacher, and his social standing was low. The Greek custom of having the boy accompanied to and from school by a slave that was otherwise incapacitated by age or physical disability soon came to be imitated by the Romans. When a special building was employed for the school, it was usually a mere booth or veranda, and the pupils sat on the floor or upon stones.

\textbf{Grammar Schools.} - The `grammar' school grew out of the increasing literary work of the \emph{Iudus.} But, while offering a more advanced course, it would seem to belong in part at least to the elementary stage of education, especially as its work was never sharply divided from that of the \emph{Iudus.} The young Roman might attend both a Greek and a Latin grammar school, but, in case he did, usually went first to the former. The curriculum in each consisted, according to Quintilian, of `the art of speaking correctly' and `the interpretation of the poets,' or, in other words, of a training in grammar and literature. `Grammar' may, however, have included some knowledge of philology and derivations, as well as drill on the parts of speech, inflections, syntax, and prosody, and practice in composition and paragraphing. The literary training was obtained by writing paraphrases of the best authors, textual and literary criticism, commentaries, and exercises in diction and verse-writing. Some other studies, like arithmetic, geometry, astronomy, geography, and music may also have been added in time, from the suggestions of Plato, but the Romans naturally gave them a practical bearing. Some gymnastics, mostly for military training, were often in the course. The methods in the grammar schools were somewhat better than those of the \emph{Indus,} but the commentary of the teacher on the text was usually taken down \emph{verbatim} by the pupil. The discipline, in consequence, was not much in advance of that of the lower schools. But the accommodations for these secondary schools were decidedly superior, and the buildings not only possessed suitable seats for the pupils and teacher, but were even adorned with paintings and sculpture.

\textbf{Rhetorical Schools.} - The `rhetorical' schools were a development of work in debate that had gradually grown up in the grammar schools. The earliest of these institutions at Rome were Greek, but by the first century B. C, there had arisen a number in which Latin was used. While they afforded a legal and forensic training, and seem more professional in spirit than the grammar schools, they were by no means narrow. The orator was for the Roman the typical man of culture and education, and he was supposed not only to have been trained in eloquence and law and history, but to possess wide learning, grace, culture, and knowledge of human emotions, sound judgment, and good memory. Besides a training in oratory, these schools furnished a linguistic and literary education of some breadth. They may be considered as belonging partly to the secondary and partly to the higher stage of education. The youths were exercised first in declamation on ethical and political subjects, which would bring in fine distinctions in Roman law and ethics, and later they were given practice in three types of speeches, - deliberative, judicial, and panegyric. Attention was given to all the various factors in making a successful oration: the matter, arrangement, style, memorizing, and delivery.

\textbf{Universities.} - When the young Roman had completed his course at a rhetorical school, he might, if he were ambitious, go to the university at Athens, Alexandria, or Rhodes for a higher training. Later, a university also sprang up at Rome, and before long these institutions spread throughout the empire. The Greek influence caused a large number of these institutions to be established in the East, but some were also located in the West. The latter gave more emphasis to practical subjects. In several instances the universities found their nucleus in one of the many libraries that were started with books brought from the sacking of Greece and Asia Minor.

\textbf{Subsidization of Education.} - Thus, through the adoption of the institutions of the Greeks, Roman education became thoroughly Hellenized. Although all the types of schools spread everywhere in the empire, there was, of course, no such thing as a real school system, except as the government gradually came to subsidize all schools. This the different emperors accomplished in various ways, - by contributing to school support, paying a salary to certain teachers, or granting them exemption from taxation and military service, or offering scholarships to a given number of pupils. As a result, schools came to be established in many cases for the purpose of getting these special privileges for the teachers, rather than for promoting education. To stop these abuses, the emperor in 425 A. D. decreed that he had the sole authority to establish schools, and that a penalty would be laid upon anyone else assuming this prerogative. In this way the schools came fully into the hands of the imperial government, and the basis for the idea of public education was laid for the first time in history.

\textbf{Decay of Education.} - Before this, however, Roman education had deteriorated. With the political and moral decay that were obvious after the second century superficial A. D., it became a mere form and mark of the aristocracy. The training in oratory was continued, because it was a necessary qualification for entering the senatorial class, but it had lost its real function, since there was no longer any occasion for oratory when the emperor dominated all the government and law. It was not intended to furnish a training of any value in life, and the careful literary preparation was more and more shirked. While the grammarians and rhetoricians were still held in high esteem, they contented themselves with mere display, and wandered from town to town more for the purpose of entertaining than of teaching. Glittering phrases, epigrams, and other artificialities took the place of instruction and argument.

\textbf{Influence of Roman Education.} - But the Roman education and civilization had left their impress upon the world. This was accomplished by the practical nature of the Romans, and by their ability to make abstract ideals concrete and embody them in institutions that have been useful to civilization and progress. Through them was created the idea of a universal empire, which has been influential throughout the world's history. Similarly, the concept of law originating with the Greek philosophers became in the hands of the Romans the great system of principles that underlies and guides all our present civilization. And it was the Roman genius for organization that institutionalized a despised religious sect and expanded it into the position of the greatest world religion. If Judaism furnished the world with exalted religious ideals, and if from Hellenism came striking intellectual and aesthetic concepts, the institutions for realizing these ideals originated with Rome.

\textbf{Supplementary Reading}

Graves, \emph{Before the Middle Ages} (Macmillan, 1909), chap. XIII; Monroe, \emph{Text-book} (Macmillan, 1905), chap. IV. Interesting brief monographs on the subject are Clarke, G., \emph{Education of Children at Rome} (Macmillan, 1896), and Wilkins, A. S., \emph{Roman Education,} (Cambridge University Press, 1905). See also the treatment in Laurie, \emph{Pre-Christian Education} (Longmans, Green, 1900), pp.~319436.

\bibliography{book.bib,packages.bib}

\end{document}
